\section{Pipeline}
\label{section:pipeline}

Figure~\ref{fig:graphicalAbstract} shows how the agents connect. For this stage each agent has a first, deliberately simple version, so that we can measure every part before combining them (Table~\ref{tab:pipeline}). The agent orchestration and the clinical report are the subject of the next stage.

\begin{table}[H]
\caption{First version of each agent.}
\label{tab:pipeline}
\centering
\small
\setstretch{1.1}
\begin{tabularx}{\textwidth}{@{}>{\raggedright\arraybackslash}p{3.6cm}>{\raggedright\arraybackslash}X>{\raggedright\arraybackslash}p{3.6cm}@{}}
\toprule
\textbf{Agent} & \textbf{First version} & \textbf{Code} \\
\midrule
Signal Monitoring & band-pass filter vs.\ convolutional denoising autoencoder & \texttt{denoise.py} \\
Arrhythmia Detection & gradient-boosted trees on RR and shape features vs.\ a 1-D CNN & \texttt{classify.py} \\
Physiological Prediction & last value (persistence) vs.\ a linear autoregressive model & \texttt{forecast.py} \\
Twin Calibration & McSharry model fitted to the patient's first 5 minutes & \texttt{twin.py} \\
Clinical Recommendation & planned for Seminar~6 & -- \\
\bottomrule
\end{tabularx}
\end{table}
